\scanpage{51}
Sinon, considérons l'un des \inserted{$\unsure{\mathcal{R}_1|R_1}$-}onglets, \deleted{\unsure{dis}} -- il y en a au moins deux. \inserted{Aussi l'un de} ces onglets \unsure{sont|est|seront} en corr. $1$-$1$. avec un ens. de \unsure{(}segments $\Sigma_2$ de $S$, dont chacun contient un composant au moins de $\deleted{S_1}T_1$, et dont les extrémités sont (\deleted{\unsure{di|dis}}distinctes et) dans $S_1$, et qui \unsure{sont} mutuellement disjoints. Soit $T_2$

\diagram{$\begin{array}{ccccc}\Sigma_1&\Sigma_2&\Sigma_3&-&\Sigma_i\\[3pt]T_1&T_2&T_3&&T_i\end{array}$}

$\Sigma_i$ est un ens. de segments \inserted{($R$-saturés)} mutuellement disjoints, $\operatorname{card}\Sigma_i\geq 2$, \unsure{$\cdot$|.} de réunion $T_i'$ \inserted{(\uline{NB} les \unsure{segm.|segmts} composants de $T_i$ sont donc des réunions partielles de $\Sigma_i$, donc saturés)}

\reversemarginpar\margin{\deleted{$i\unsure{\geq|\to}$}} \[T_i=T_i'\cup T_{\unsure{i-1|i}}\] \margin{$i\geq 2$ \deleted{\unsure{donc}}}\noindent\llap{\rule[-2\baselineskip]{.4pt}{3\baselineskip}\hspace{.4em}}\deleted{Tout élément $A\in\Sigma_{i-1}$ est contenu dans un $\deleted{\unsure{A}}B\in\Sigma_i$, \unsure{les|tout}}

Tout $B\in\Sigma_i$ contient au moins un \deleted{$A\in\Sigma_{i-1}$}\inserted{\uline{segment composant}} \deleted{\unsure{de}} de $T_{i-1}$, et $\partial B\subset S\setminus T_{i-1}$ \deleted{$\unsure{\partial B|B}$}\normalmarginpar

\noindent\begin{minipage}[t]{.46\linewidth}
$\Sigma_1$ ens de segments \unsure{mut\textsuperscript{t}|mut\textsuperscript{lt}} disj. de $S=S_0$, \deleted{de cardinal} \inserted{$\operatorname{card}(\Sigma_1)\geq 2$} \[T_1=\bigcup_{A\in\Sigma_1} A,\quad S_1=S\setminus T_1\] \[\deleted{\operatorname{card}S_1\geq\unsure{3|2}}\] les segm. comp. de $S_1$ dans $S$ au \unsure{nb|nbre} au moins $2$
\end{minipage}\hfill\begin{minipage}[t]{.51\linewidth}
$\Sigma'_2$ ens de segm. mut. disjoints \uline{de $S_{\unsure{2|1}}$}, \uline{non} \uline{contenus} dans un des segments composants \[\begin{aligned}\operatorname{card}\Sigma_{\unsure{1|2}}&\geq 2,\\ T_2&=T_1\cup\mathop{\deleted{\Sigma}\bigcup}\limits_{A\in\Sigma'_2} A\\ \deleted{\unsure{\Sigma}}S_2&=S\setminus T_2\\ \operatorname{card}\unsure{\pi_0|\Pi_0}(S_2)&\ne 1\end{aligned}\]
\end{minipage}
% archive sheet number: 50

